\section*{ID6510: Compact Analytic Model: Dipole + Axial Vortex + Wave Modes: γ}
\paragraph{Metadata.}
PiID: 8581538420568 | RTEID: RTE:P25279.7757:T3.9542 | UUID: 57494054-15a0-512f-b79b-a23347dd5109

\paragraph{Canonical Statement.}
**Axial vortex (localized pull around a pole — can bias one arm).** Model a concentrated vortex centered at the north pole (a similar one can be put on the south pole or both): A classic streamfunction for a point vortex at colatitude \textbackslash{}(\textbackslash{}theta\_0\textbackslash{}) on the sphere (angular distance \textbackslash{}(\textbackslash{}gamma\textbackslash{})) is \textbackslash{}[ \textbackslash{}Psi\_\{\textbackslash{}rm vortex\}(\textbackslash{}gamma)=\textbackslash{}frac\{\textbackslash{}Gamma\}\{4\textbackslash{}pi\}\textbackslash{}ln\textbackslash{}big(1-\textbackslash{}cos\textbackslash{}gamma\textbackslash{}big), \textbackslash{}] where for north-pole-centered vortex \textbackslash{}(\textbackslash{}gamma=\textbackslash{}theta\textbackslash{}). This gives a tangential circulation decaying away from the pole and is useful to represent stronger pull in one hemisphere.

\paragraph{Canonical Equation.}
\[
\gamma=\theta
\]

\paragraph{Dictionary.}
Axial vortex (localized pull around a pole — can bias one arm) (γ)

\paragraph{Encyclopedia.}
Compact Analytic Model: Dipole + Axial Vortex + Wave Modes: γ is an algebraic definition, written γ=θ. It is an algebraic definition expressing the quantity directly in terms of the variables on its right-hand side. It is built from θ, defining γ. Within the Trawin Zero Point registry it serves as one element of the framework's mathematical narrative.
